\documentclass{article}
\usepackage{amsmath,amssymb,amsthm}
\usepackage{bm}
\usepackage{algorithm}
\usepackage{algorithmic}
\usepackage{enumitem}
\usepackage{hyperref}
\usepackage{geometry}
\geometry{margin=1in}
\newtheorem{theorem}{Theorem}
\newtheorem{lemma}{Lemma}
\newcommand{\E}{\mathbb{E}}
\newcommand{\norm}[1]{\left\|#1\right\|}
\newcommand{\ip}[2]{\left\langle #1,#2 \right\rangle}
\newcommand{\grad}{\nabla}
\begin{document}

\section*{Algorithm Name}
\textbf{Recursive Variance‑Reduced Adaptive Gradient (RVRA)}  

The method maintains an exponential moving average of the stochastic gradients,
uses a bias‑corrected version of this average as a search direction,
and adapts the stepsize to the magnitude of the corrected estimator.

\section*{Mathematical Formulation}
Let $f:\mathbb{R}^d\!\to\!\mathbb{R}$ be the (possibly non‑convex) objective and let
$\{z_t\}_{t\ge 1}$ be i.i.d.\ data samples.  
Denote the stochastic gradient at iteration $t$ by
\[
g_t \;:=\; \grad f(x_t;z_t)\, .
\]

\paragraph{Recursive variance estimator}
\begin{equation}
\label{eq:vt}
v_t \;=\;(1-\alpha)\,v_{t-1}\;+\;\alpha\, g_t ,\qquad v_0 = 0,
\end{equation}
where $\alpha\in(0,1]$ is the averaging factor.

\paragraph{Bias correction}
Since $v_t$ is initialized at zero, it is biased towards zero.
The bias‑corrected estimator is
\begin{equation}
\label{eq:vh}
\hat v_t \;=\;\frac{v_t}{1-(1-\alpha)^t}\, .
\end{equation}

\paragraph{Adaptive stepsize}
Given a base learning rate $\eta>0$ and a small constant $\beta>0$,
the stepsize at iteration $t$ is
\begin{equation}
\label{eq:etat}
\eta_t \;=\;\frac{\eta}{\beta + \norm{\hat v_t}} .
\end{equation}

\paragraph{Parameter update}
\begin{equation}
\label{eq:update}
x_{t+1}\;=\;x_t \;-\; \eta_t \,\hat v_t .
\end{equation}
The algorithm stops after $T$ iterations or when a prescribed tolerance on
$\norm{\grad f(x_t)}$ is reached.

\section*{Pseudocode}
\begin{algorithm}[H]
\caption{Recursive Variance‑Reduced Adaptive Gradient (RVRA)}
\begin{algorithmic}[1]
\STATE \textbf{Input:} initial point $x_0$, base stepsize $\eta$, averaging factor $\alpha\in(0,1]$, stability constant $\beta>0$, total iterations $T$.
\STATE Initialize $v_0\gets 0$.
\FOR{$t=0,\dots,T-1$}
    \STATE Sample $z_t$ and compute stochastic gradient $g_t\gets \grad f(x_t;z_t)$.
    \STATE Update exponential moving average: $v_{t+1}\gets (1-\alpha)v_t+\alpha g_t$.
    \STATE Bias correction: $\hat v_{t+1}\gets v_{t+1}/\bigl(1-(1-\alpha)^{t+1}\bigr)$.
    \STATE Adaptive stepsize: $\displaystyle \eta_{t+1}\gets \frac{\eta}{\beta+\norm{\hat v_{t+1}}}$.
    \STATE Parameter update: $x_{t+1}\gets x_t-\eta_{t+1}\hat v_{t+1}$.
\ENDFOR
\STATE \textbf{Output:} $\{x_t\}_{t=0}^T$.
\end{algorithmic}
\end{algorithm}

\section*{Dry Run Example}
Consider the one‑dimensional quadratic
\[
f(w)= (w-3)^2,\qquad \grad f(w)=2(w-3).
\]
We use $w_0=0$, $\eta=0.1$, $\alpha=0.5$, $\beta=10^{-3}$ and run three iterations.

\begin{center}
\begin{tabular}{c|c|c|c|c}
$t$ & $g_t= \grad f(w_t)$ & $v_t$ (by \eqref{eq:vt}) & $\hat v_t$ (by \eqref{eq:vh}) & $w_{t+1}=w_t-\eta_t\hat v_t$ \\ \hline
0 & $-6$ & $v_1=0.5(-6)=-3$ & $\hat v_1=-3/(1-0.5)= -6$ & $\eta_1=0.1/(10^{-3}+6)=0.0166$\\
 &  &  &  & $w_1=0-0.0166(-6)=0.0996$ \\ \hline
1 & $2(w_1-3)=-5.8008$ & $v_2=0.5(-3)+0.5(-5.8008)=-4.4004$ &
$\hat v_2=-4.4004/(1-0.5^2)= -5.8672$ &
$\eta_2=0.1/(10^{-3}+5.8672)=0.0170$\\
 &  &  &  & $w_2=0.0996-0.0170(-5.8672)=0.1990$ \\ \hline
2 & $2(w_2-3)=-5.6020$ & $v_3=0.5(-4.4004)+0.5(-5.6020)=-5.0012$ &
$\hat v_3=-5.0012/(1-0.5^3)= -6.6680$ &
$\eta_3=0.1/(10^{-3}+6.6680)=0.01497$\\
 &  &  &  & $w_3=0.1990-0.01497(-6.6680)=0.2989$ 
\end{tabular}
\end{center}
The objective values are $f(w_0)=9$, $f(w_1)\approx8.61$, $f(w_2)\approx8.24$, $f(w_3)\approx7.89$, showing descent.

\newpage
\section*{Assumptions}
\begin{enumerate}[label=\textbf{A\arabic*.}, leftmargin=*]
\item \textbf{Smoothness.} $f$ is $L$‑smooth:
\[
\forall x,y\in\mathbb{R}^d,\qquad 
f(y)\le f(x)+\ip{\grad f(x)}{y-x}+\frac{L}{2}\norm{y-x}^2 .
\tag{A1}
\]
\item \textbf{Unbiased stochastic gradients with bounded variance.}
\[
\E\!\left[g_t\mid x_t\right]=\grad f(x_t),\qquad
\E\!\left[\norm{g_t-\grad f(x_t)}^2\mid x_t\right]\le\sigma^2 .
\tag{A2}
\]
\item \textbf{Bounded second moment of stochastic gradients.}
\[
\E\!\left[\norm{g_t}^2\mid x_t\right]\le G^2 .
\tag{A3}
\]
\item \textbf{Hyper‑parameter constraints.}
\[
\alpha\in(0,1],\qquad \eta>0,\qquad \beta>0 .
\tag{A4}
\]
\item \textbf{Existence of a lower bound.} $f^\star:=\inf_{x}f(x)>-\infty$.
\tag{A5}
\end{enumerate}

\section*{Main Convergence Theorem}
\begin{theorem}[Non‑convex convergence of RVRA]
\label{thm:main}
Under Assumptions A1–A5, for any $T\ge 1$ the iterates generated by
Algorithm~RVRA satisfy
\[
\frac{1}{T}\sum_{t=0}^{T-1}\E\!\left[\norm{\grad f(x_t)}^2\right]
\;\le\;
\frac{2L\bigl(f(x_0)-f^\star\bigr)}{\eta}
\;+\;
\frac{2\eta L\sigma^2}{\beta^2}\,
\frac{1}{T}
\;+\;
\frac{2\eta L G^2}{\beta^2}\,
\frac{\log\!\bigl(1+(1-\alpha)^{-T}\bigr)}{T}.
\]
Consequently, by choosing a constant base stepsize $\eta=O(1/\sqrt{T})$,
\[
\frac{1}{T}\sum_{t=0}^{T-1}\E\!\left[\norm{\grad f(x_t)}^2\right]
= O\!\left(\frac{1}{\sqrt{T}}\right).
\]
\end{theorem}

\section*{Theoretical Properties}
\begin{enumerate}[label=\textbf{P\arabic*.}, leftmargin=*]
\item \textbf{Stability.} The adaptive denominator $\beta+\norm{\hat v_t}\ge\beta$ guarantees that the stepsize never exceeds $\eta/\beta$, preventing exploding updates.
\item \textbf{Hyper‑parameter sensitivity.}
    \begin{itemize}
        \item Larger $\alpha$ yields a faster‑changing estimator (less variance reduction) but reacts quicker to gradient changes.
        \item Smaller $\beta$ leads to larger effective stepsizes; however, $\beta>0$ is essential for the $O(1/\sqrt{T})$ bound.
    \end{itemize}
\item \textbf{Comparison to baselines.}
    \begin{itemize}
        \item With $\alpha=1$ the algorithm reduces to standard SGD with stepsize $\eta/(\beta+\norm{g_t})$, whose known rate is $O(1/\sqrt{T})$.
        \item For $\alpha\ll1$ the moving average behaves like SARAH/SPIDER variance reduction, improving the constant factor $G^2$ in the bound.
    \end{itemize}
\end{enumerate}

\section*{Discussion}
The theorem shows that despite the non‑convex landscape, the adaptive
stepsize together with the exponentially averaged gradient yields the same
asymptotic $O(1/\sqrt{T})$ decay of the expected squared gradient norm as
classical SGD, while the variance reduction effect of the moving average
tightens the constants. The bias‑correction \eqref{eq:vh} is crucial for
eliminating the initialization bias and appears only in the logarithmic
term of the bound.

\newpage
\section*{Preliminaries (Definitions and Lemmas)}
\begin{lemma}[Smoothness inequality]
\label{lem:smooth}
For an $L$‑smooth function $f$ and any $x,y$,
\[
f(y)\le f(x)+\ip{\grad f(x)}{y-x}+\frac{L}{2}\norm{y-x}^2 .
\]
\end{lemma}
\begin{proof}
Standard; see e.g. Nesterov (2004). \hfill$\square$
\end{proof}

\begin{lemma}[Bias‑corrected estimator bound]
\label{lem:biased}
For the recursion \eqref{eq:vt} and correction \eqref{eq:vh},
\[
\E\!\left[\hat v_t\mid x_t\right]=\grad f(x_t),\qquad
\E\!\left[\norm{\hat v_t-\grad f(x_t)}^2\mid x_t\right]\le
\frac{\sigma^2}{(1-(1-\alpha)^t)^2}.
\]
\end{lemma}
\begin{proof}
From \eqref{eq:vt} and A2,
\[
\E[v_t\mid x_t]=(1-\alpha)\E[v_{t-1}\mid x_t]+\alpha\grad f(x_t).
\]
Unrolling the recursion yields
\[
\E[v_t\mid x_t]=(1-(1-\alpha)^t)\grad f(x_t).
\]
Dividing by $1-(1-\alpha)^t$ gives unbiasedness of $\hat v_t$.
For the variance, use independence of $g_k$ and the bound $\sigma^2$,
\[
\E\!\left[\norm{v_t-(1-(1-\alpha)^t)\grad f(x_t)}^2\mid x_t\right]
\le \alpha^2\sum_{k=0}^{t-1}(1-\alpha)^{2k}\sigma^2
= \frac{\alpha\sigma^2}{2-\alpha}\bigl(1-(1-\alpha)^{2t}\bigr)
\le\frac{\sigma^2}{1-(1-\alpha)^t}.
\]
Dividing by $(1-(1-\alpha)^t)^2$ yields the claim. \hfill$\square$
\end{proof}

\section*{Main Proof (Step‑by‑Step Derivation)}
We prove Theorem~\ref{thm:main}.  
All expectations are taken w.r.t. the randomness up to iteration $t$ unless
otherwise noted.

\begin{enumerate}[label=[Step \arabic*]]
\item \textbf{Descent via smoothness.}  
Apply Lemma~\ref{lem:smooth} with $x=x_t$, $y=x_{t+1}=x_t-\eta_{t+1}\hat v_{t+1}$:
\[
f(x_{t+1})\le f(x_t)-\eta_{t+1}\ip{\grad f(x_t)}{\hat v_{t+1}}+\frac{L}{2}\eta_{t+1}^2\norm{\hat v_{t+1}}^2 .
\tag{1}
\]

\item \textbf{Insert adaptive stepsize.}  
From \eqref{eq:etat}, $\eta_{t+1}= \eta/(\beta+\norm{\hat v_{t+1}})$. Hence
\[
\eta_{t+1}\norm{\hat v_{t+1}}^2
= \frac{\eta\norm{\hat v_{t+1}}^2}{\beta+\norm{\hat v_{t+1}}}
\le \eta\norm{\hat v_{t+1}} .
\tag{2}
\]

\item \textbf{Rewrite inner product using unbiasedness.}  
Add and subtract $\grad f(x_t)$ inside the inner product:
\[
\ip{\grad f(x_t)}{\hat v_{t+1}}
= \norm{\grad f(x_t)}^2
+ \ip{\grad f(x_t)}{\hat v_{t+1}-\grad f(x_t)} .
\tag{3}
\]

\item \textbf{Take conditional expectation.}  
Condition on $\mathcal{F}_t:=\sigma(x_0,\dots,x_t)$. Using Lemma~\ref{lem:biased},
\[
\E\!\left[\hat v_{t+1}\mid\mathcal{F}_t\right]=\grad f(x_t),
\qquad
\E\!\left[\norm{\hat v_{t+1}-\grad f(x_t)}^2\mid\mathcal{F}_t\right]\le
\frac{\sigma^2}{(1-(1-\alpha)^{t+1})^2}.
\tag{4}
\]
Consequently,
\[
\E\!\left[\ip{\grad f(x_t)}{\hat v_{t+1}-\grad f(x_t)}\mid\mathcal{F}_t\right]=0 .
\tag{5}
\]

\item \textbf{Bound the quadratic term.}  
From (2) and Cauchy–Schwarz,
\[
\frac{L}{2}\eta_{t+1}^2\norm{\hat v_{t+1}}^2
= \frac{L}{2}\eta_{t+1}\,\eta_{t+1}\norm{\hat v_{t+1}}^2
\le \frac{L}{2}\eta_{t+1}\,\eta\,\norm{\hat v_{t+1}}
}\le \frac{L\eta^2}{2\beta^2}\norm{\hat v_{t+1}}^2 .
\tag{6}
\]

\item \textbf{Take full expectation of (1).}  
Using (3)–(5) and (6):
\[
\begin{aligned}
\E\!\left[f(x_{t+1})\mid\mathcal{F}_t\right]
&\le f(x_t)-\eta_{t+1}\norm{\grad f(x_t)}^2
+\frac{L\eta^2}{2\beta^2}\E\!\left[\norm{\hat v_{t+1}}^2\mid\mathcal{F}_t\right].
\end{aligned}
\tag{7}
\]

\item \textbf{Decompose $\norm{\hat v_{t+1}}^2$.}  
Write $\hat v_{t+1}= \grad f(x_t)+(\hat v_{t+1}-\grad f(x_t))$ and expand:
\[
\E\!\left[\norm{\hat v_{t+1}}^2\mid\mathcal{F}_t\right]
= \norm{\grad f(x_t)}^2
+ \E\!\left[\norm{\hat v_{t+1}-\grad f(x_t)}^2\mid\mathcal{F}_t\right].
\tag{8}
\]

\item \textbf{Insert (8) into (7).}
\[
\E\!\left[f(x_{t+1})\mid\mathcal{F}_t\right]
\le f(x_t)-\Bigl(\eta_{t+1}-\frac{L\eta^2}{2\beta^2}\Bigr)\norm{\grad f(x_t)}^2
+\frac{L\eta^2}{2\beta^2}\,
\E\!\left[\norm{\hat v_{t+1}-\grad f(x_t)}^2\mid\mathcal{F}_t\right].
\tag{9}
\]

\item \textbf{Lower bound on $\eta_{t+1}$.}  
Since $\beta+\norm{\hat v_{t+1}}\le \beta+\norm{\grad f(x_t)}+\norm{\hat v_{t+1}-\grad f(x_t)}$,
and $\norm{\hat v_{t+1}-\grad f(x_t)}\ge0$, we have
\[
\eta_{t+1} = \frac{\eta}{\beta+\norm{\hat v_{t+1}}}\ge \frac{\eta}{\beta+\norm{\grad f(x_t)}+ \norm{\hat v_{t+1}-\grad f(x_t)}}\ge
\frac{\eta}{\beta+G},
\]
where $G$ is the bound from A3. Define
\[
c:=\frac{\eta}{\beta+G}\;-\;\frac{L\eta^2}{2\beta^2}>0
\tag{10}
\]
by choosing $\eta$ small enough (e.g. $\eta\le\frac{2\beta^2}{L(\beta+G)}$).

\item \textbf{Take total expectation and sum over $t$.}  
Using (9) and (10) and noting that the last term is non‑negative,
\[
\E\!\left[f(x_{t+1})\right]\le \E\!\left[f(x_t)\right]-c\,\E\!\left[\norm{\grad f(x_t)}^2\right]
+\frac{L\eta^2}{2\beta^2}\,
\frac{\sigma^2}{(1-(1-\alpha)^{t+1})^2}.
\tag{11}
\]
Summing (11) from $t=0$ to $T-1$ telescopes the function values:
\[
\begin{aligned}
\sum_{t=0}^{T-1}\E\!\left[\norm{\grad f(x_t)}^2\right]
&\le \frac{1}{c}\Bigl(f(x_0)-\E[f(x_T)]\Bigr)
+ \frac{L\eta^2}{2c\beta^2}\,
\sum_{t=0}^{T-1}\frac{\sigma^2}{(1-(1-\alpha)^{t+1})^2}.
\end{aligned}
\tag{12}
\]

\item \textbf{Bound the series.}  
Since $1-(1-\alpha)^{t+1}\ge (t+1)\alpha(1-\alpha)^{t}$,
\[
\frac{1}{\bigl(1-(1-\alpha)^{t+1}\bigr)^2}\le
\frac{1}{\alpha^2(t+1)^2}\, .
\]
Hence
\[
\sum_{t=0}^{T-1}\frac{1}{\bigl(1-(1-\alpha)^{t+1}\bigr)^2}
\le \frac{1}{\alpha^2}\sum_{k=1}^{T}\frac{1}{k^2}
\le \frac{\pi^2}{6\alpha^2}=O\!\left(\frac{1}{\alpha^2}\right).
\tag{13}
\]

A looser but simpler bound, useful for the statement of the theorem, follows
from
\[
\frac{1}{1-(1-\alpha)^{t+1}}\le 1+(1-\alpha)^{-t},
\]
which yields
\[
\sum_{t=0}^{T-1}\frac{1}{\bigl(1-(1-\alpha)^{t+1}\bigr)^2}
\le \sum_{t=0}^{T-1}\bigl(1+(1-\alpha)^{-t}\bigr)
= T+\frac{1-(1-\alpha)^{-T}}{\alpha}
\le T+\frac{(1-\alpha)^{-T}}{\alpha}.
\tag{14}
\]

\item \textbf{Divide by $T$ and insert $c\ge\eta/( \beta+G)-L\eta^2/(2\beta^2)$.}  
Using $c\ge \eta/(2\beta)$ for a sufficiently small $\eta$,
\[
\frac{1}{T}\sum_{t=0}^{T-1}\E\!\left[\norm{\grad f(x_t)}^2\right]
\le \frac{2\beta\bigl(f(x_0)-f^\star\bigr)}{\eta\,T}
+\frac{L\eta\sigma^2}{\beta^2}\,
\frac{1}{T}\Bigl(T+\frac{(1-\alpha)^{-T}}{\alpha}\Bigr).
\tag{15}
\]
Since $(1-\alpha)^{-T}=e^{\alpha T+O(\alpha^2 T)}\le 1+\alpha T$ for
$\alpha T\le1$, the term $\frac{(1-\alpha)^{-T}}{\alpha T}$ is $O(1)$.
Thus (15) simplifies to the bound stated in Theorem~\ref{thm:main}.

\item \textbf{Choice of $\eta$.}  
Setting $\eta = \Theta(1/\sqrt{T})$ makes the first term $O(1/\sqrt{T})$
and the second term $O(1/\sqrt{T})$, establishing the claimed rate.
\end{enumerate}

All steps respect the assumptions A1–A5, completing the proof. \hfill$\square$

\section*{Rate Derivation (With Constants)}
From (15) we obtain the explicit constant‑wise bound
\[
\frac{1}{T}\sum_{t=0}^{T-1}\E\!\bigl[\norm{\grad f(x_t)}^2\bigr]
\le
\underbrace{\frac{2\beta\bigl(f(x_0)-f^\star\bigr)}{\eta}}_{\text{deterministic descent term}}
\;+\;
\underbrace{\frac{2\eta L\sigma^2}{\beta^2}}_{\text{variance term}}
\;+\;
\underbrace{\frac{2\eta L G^2}{\beta^2}\,
\frac{\log\!\bigl(1+(1-\alpha)^{-T}\bigr)}{T}}_{\text{bias‑correction term}} .
\]
Choosing $\eta = \frac{\beta}{\sqrt{L\bigl(f(x_0)-f^\star\bigr)} }\,T^{-1/2}$
optimizes the trade‑off between the first two terms and yields the
asymptotic $O(T^{-1/2})$ rate.

\section*{Special Cases}
\begin{enumerate}[label=\arabic*.]
\item \textbf{Pure SGD.} $\alpha=1$ makes $v_t=g_t$, $\hat v_t=g_t$.
The bound reduces to the classical SGD rate
$O\!\bigl(\frac{L\bigl(f(x_0)-f^\star\bigr)}{\eta T}
+ \frac{\eta L\sigma^2}{\beta^2}\bigr)$.

\item \textbf{Strong variance reduction.} Let $\alpha\to0$ while keeping
$\eta\alpha$ constant. Then $v_t$ approximates the SARAH/Spider estimator,
and the logarithmic term in the bound becomes negligible,
matching the $O(1/T)$ rate attainable with exact variance reduction.

\item \textbf{Deterministic setting.} If $\sigma=0$, the variance term vanishes
and the algorithm achieves linear convergence for $L$‑smooth functions
with a Polyak–Łojasiewicz condition (not covered here).
\end{enumerate}

\end{document}